\section*{Resumen}
La agricultura cubana enfrenta una crisis multiplicativa caracterizada por la escasez de materias primas, energía, agua y degradación de suelos, exacerbada por el cambio climático y el bloqueo económico. La infraestructura de riego se encuentra colapsada, con tecnologías obsoletas y ausencia de precisión agrícola, mientras la gestión institucional resulta deficiente por programación de riego sin balances hídricos y falta de personal capacitado. La intersección con la crisis energética—apagones programados, falta de combustible y piezas de repuesto—paraliza los sistemas de bombeo, incluso aquellos más eficientes. Además, la salinización progresiva de los suelos por sobreexplotación sin drenaje adecuado genera estrés salino que reduce la productividad.
Este trabajo propone la asimilación tecnológica del Internet de las Cosas (IoT) como enfoque para mitigar la mala gestión de recursos hídricos en el contexto agrícola cubano. Se describirá el entorno actual de tecnologías de transmisión de datos en tiempo real, se analizarán los artefactos necesarios para generar un sistema de métricas para el control de recursos, y se ejecutará la asimilación de la tecnología en condiciones de restricción económica y energética. El objeto de estudio se centra en las herramientas de transmisión de datos en entornos IoT aplicadas a la agricultura de precisión, buscando contribuir a la soberanía alimentaria y seguridad nacional establecidas en las políticas gubernamentales hasta 2030.


%\begin{description}
%	\item[Palabras claves:]{la lista de palabras separadas por comas (,)}
%\end{description}
%\end Resumen


